\documentclass[12pt,a4paper]{article}
\usepackage[utf8]{inputenc}
\usepackage[T1]{fontenc}
\usepackage[polish]{babel}
\usepackage{geometry}
\usepackage{graphicx}
\usepackage{titlesec}
\usepackage{indentfirst}
\usepackage{listings}
\usepackage{xcolor}
\usepackage{setspace}
\usepackage{hyperref} % Hyperref warto ładować na końcu preambuły

\geometry{
	a4paper,
	total={170mm,257mm},
	left=20mm,
	top=20mm,
}
\onehalfspacing

\title{\textbf{Dokumentacja Projektowa}\\Systemy Przetwarzania w Chmurze}
\author{Zespół Projektowy}
\date{\today}
\newpage

\begin{document}
	
	\maketitle
	\tableofcontents
	\newpage
	
	\section{Wprowadzenie}
	Niniejsza dokumentacja stanowi podsumowanie prac nad projektem \textit{GutCloud} realizowanym w ramach zajęć projektowych z przedmiotu "Systemy Przetwarzania w Chmurze". Głównym celem było zaprojektowanie, zaimplementowanie oraz wdrożenie systemu informatycznego opartego na modelu przetwarzania w chmurze.
	
	Stworzony system to aplikacja webowa służąca do zarządzania plikami użytkowników. Jest to odpowiedź na rosnące zapotrzebowanie na bezpieczne, skalowalne i godne zaufania repozytoria danych. Tradycyjne metody przechowywania danych na dyskach lokalnych stają się niewystarczające dla użytkowników korzystających z wielu urządzeń jednocześnie i oczekujących dostępu do swoich zasobów. Nasza aplikacja rozwiązuje ten problem, wykorzystując infrastrukturę chmurową, gwarantującą wysoką dostępność i trwałość.
	
	System umożliwia użytkownikom bezpieczną metodę logowania oferowaną przez dostawcę usługi chmurowej oraz wykonywanie podstawowych operacji na plikach. Zaimplementowane mechanizmy umożliwiają wygodne przesyłanie całych struktur katalogów. Zastosowane zostały rozwiązania z zakresu Smart City, poprzez integrację systemu z algorytmami sztucznej inteligencji. Aplikacja wykorzystuje model językowy do generowania inteligentnych podsumowań zawartości plików tekstowych. Dzięki temu użytkownik nie musi otwierać każdego dokumentu, aby poznać jego treść.
	
	Architektura rozwiązania opiera się na sprawdzonym wzorcu Model-View-Controller zaimplementowanym przy użyciu języka Python i frameworka Flask. Wybór tej technologii był podyktowany faktem, iż Python jest obecnie jednym z najpopularniejszych języków w świecie analizy danych i technologii chmurowych, co ułatwiło integrację z bibliotekami SDK Azure oraz API sztucznej inteligencji. Warstwa danych została całkowicie odseparowana od warstwy aplikacji poprzez wykorzystanie usług Azure Blob Storage. Taka separacja zapewnia skalowalność, bez ingerencji w warstwę danych.
	
	Proces realizacji projektu obejmował kilka etapów: od analizy wymagań, przez projektowanie architektury, implementację poszczególnych modułów, aż po wdrożenie. Każdy z tych etapów pozwolił zespołowi na zdobycie cennych umiejętności w zakresie administrowania chmurą, programowania aplikacji webowych oraz pracy grupowej z wykorzystaniem systemów kontroli wersji.
	
	W dalszej części dokumentacji przedstawiono szczegółowy opis techniczny, analizę wykorzystanych technologii i usług chmurowych oraz instrukcję obsługi dla użytkownika końcowego.
	
	\newpage
	\section{Prezentacja chmury i dostawcy}
	Jako dostawcę usług chmurowych wybrano \textbf{Microsoft Azure}. Jest to jedna z czołowych platform chmurowych na rynku, a decyzja o jej wyborze podyktowana była kilkoma czynnikami:
	\begin{enumerate}
		\item \textbf Azure oferuje stabilne i sprawdzone usługi magazynowania danych, które są fundamentem naszego projektu.
		\item \textbf Dostępność doskonałych bibliotek SDK dla języka Python znacznie przyspieszyła proces rozwoju aplikacji.
		\item \textbf{Dokumentacja}: Microsoft udostępnia obszerną dokumentację techniczną oraz przykłady kodu, co było ogromną pomocą podczas nauki i rozwiązywania problemów.
		\item \textbf{Kredyty studenckie/darmowe zasoby}: Dostępność darmowych tierów usług dla studentów i deweloperów.
	\end{enumerate}
	
	Najważniejsze funkcje i usługi Azure wykorzystane w projekcie:
	\begin{itemize}
		\item \textbf{Azure Storage Account}: Kontener dla wszystkich usług danych. W jego ramach wykorzystujemy:
		\begin{itemize}
			\item \textbf{Blob Storage}: Usługa do przechowywania ogromnych ilości danych niestrukturalnych. Blob Storage jest idealny do udostępniania plików bezpośrednio w przeglądarce, przechowywania plików do dostępu rozproszonego, strumieniowania wideo i audio oraz przechowywania danych do kopii zapasowych i przywracania, odzyskiwania po awarii i archiwizacji. W naszym projekcie użyliśmy go do trwałego zapisu plików użytkowników.
			\item \textbf{Table Storage}: Usługa przechowująca ustrukturyzowane dane NoSQL w chmurze, zapewniająca magazyn kluczy/atrybutów z bezschematowym projektem. Jest to szybkie i tanie rozwiązanie dla danych, które nie wymagają złożonych złączeń SQL. W projekcie posłużyła jako baza danych użytkowników (credentials, flagi admina).
		\end{itemize}
	\end{itemize}
	
	\newpage
	\section{Prezentacja techniczna aplikacji}
	
	\subsection{Wymagania}
	System został zaprojektowany w oparciu o następujące wymagania:
	\begin{itemize}
		\item \textbf{Funkcjonalne}:
		\begin{itemize}
			\item Logowanie użytkowników.
			\item Przesyłanie plików do chmury.
			\item Pobieranie plików na urządzenie lokalne.
			\item Usuwanie plików (pojedynczo i masowo).
			\item Obsługa plików ZIP (automatyczne rozpakowywanie w tle).
			\item Integracja z AI do generowania opisów plików tekstowych.
			\item Logowanie zdarzeń systemowych (dla administratora).
		\end{itemize}
		\item \textbf{Niefunkcjonalne}:
		\begin{itemize}
			\item Bezpieczeństwo danych.
			\item Skalowalność (dzięki Azure Blob Storage).
			\item Responsywność interfejsu.
			\item Dostępność gwarantowana przez SLA dostawcy chmury.
		\end{itemize}
	\end{itemize}
	
	\subsection{Architektura i Warstwa Logiczna}
	Aplikacja działa w architekturze klient-serwer.
	\begin{itemize}
		\item \textbf{Backend}: Python 3.x + Flask. Flask pełni rolę kontrolera, obsługując routing zapytań HTTP. Wykorzystuje biblioteki \texttt{azure.storage.blob} do komunikacji z Blob Storage oraz \texttt{azure.data.tables} do obsługi Table Storage.
		\item \textbf{Frontend}: HTML5 + CSS + Bootstrap. Strony są generowane dynamicznie po stronie serwera.
		\item \textbf{Baza Danych}: Azure Table Storage (NoSQL).
	\end{itemize}
	
	Kluczowe mechanizmy logiczne:
	\begin{itemize}
		\item \textbf{Wersjonowanie}: Przy przesyłaniu pliku o istniejącej nazwie, system nie nadpisuje go, lecz dodaje timestamp do nazwy (np. \texttt{plik\_20231012.txt}).
		\item \textbf{Asynchroniczność}: Rozpakowywanie archiwów ZIP jest operacją czasochłonną, dlatego została wydzielona do osobnego wątku. Dzięki temu użytkownik natychmiast otrzymuje informację o rozpoczęciu procesu i może kontynuować pracę, podczas gdy serwer przetwarza plik w tle.
		\item \textbf{Integracja z LLM}: Wykorzystanie biblioteki \texttt{groq} do wysłania fragmentu tekstu pliku do modelu Llama 3 i odebrania wygenerowanego podsumowania.
	\end{itemize}
	
	\newpage
	\section{Aspekty Smart City i STEM}
	Aplikacja realizuje koncepcje Smart City poprzez ideę \textbf{otwartych i dostępnych danych (Open Data \& Accessibility)}. W inteligentnym mieście obywatele i urzędy muszą mieć możliwość wymiany dokumentów w sposób cyfrowy, bezpieczny i niezależny od lokalizacji. Nasz system może stanowić prototyp miejskiej platformy wymiany dokumentów, gdzie mieszkańcy składają wnioski (upload plików), a systemy AI wstępnie je kategoryzują i opisują (automatyzacja urzędów).
	
	Wymiar \textbf{STEM} (Science, Technology, Engineering, Mathematics) jest realizowany poprzez:
	\begin{itemize}
		\item \textbf{Science}: Zastosowanie najnowszych osiągnięć w dziedzinie przetwarzania języka naturalnego (LLM) do analizy treści.
		\item \textbf{Technology}: Użycie chmury obliczeniowej, która jest obecnie standardem technologicznym w przemyśle IT.
		\item \textbf{Engineering}: Zaprojektowanie i wdrożenie kompletnego systemu informatycznego z uwzględnieniem obsługi błędów, wielowątkowości i zarządzania zależnościami.
		\item \textbf{Mathematics}: Wykorzystanie kryptografii i logiki matematycznej w procesach autoryzacji i przetwarzania danych.
	\end{itemize}
	
	\newpage
	\section{Rozwiązania chmurowe}
	Kluczowe rozwiązania chmurowe zaimplementowane w projekcie:
	\begin{enumerate}
		\item \textbf{Azure Blob Storage}:
		\begin{itemize}
			\item Kontener \texttt{files} przechowuje wszystkie pliki.
			\item Ustrukturyzowane nazewnictwo blobów (\texttt{user\_id/nazwa\_pliku}) emuluje system plików, mimo płaskiej struktury blob storage. Zapewnia to separację danych między użytkownikami na poziomie aplikacji.
		\end{itemize}
		\item \textbf{Azure Table Storage}:
		\begin{itemize}
			\item Wykorzystanie \texttt{PartitionKey} i \texttt{RowKey} do szybkiego dostępu do danych użytkownika. \texttt{PartitionKey="users"} grupuje wszystkich użytkowników, a \texttt{RowKey} (login) zapewnia unikalność i szybkie wyszukiwanie.
		\end{itemize}
		\item \textbf{Groq AI Cloud}:
		\begin{itemize}
			\item Wykorzystanie mocy obliczeniowej chmury Groq do inferencji modelu językowego, co odciąża serwer aplikacji.
		\end{itemize}
	\end{enumerate}
	
	\newpage
	\section{Podręcznik użytkownika}
	\subsection{Rozpoczęcie pracy}
	\begin{enumerate}
		\item Otwórz przeglądarkę i wejdź na adres aplikacji.
		\item Jeśli nie masz konta, wybierz opcję "Zarejestruj się". Podaj unikalny login i hasło.
		\item Zaloguj się korzystając z utworzonych danych.
	\end{enumerate}
	
	\subsection{Główne operacje}
	\begin{itemize}
		\item \textbf{Przesyłanie plików}: W sekcji "Upload" wybierz plik z dysku lokalnego. Obsługiwane są wszystkie formaty. Pliki ZIP zostaną automatycznie rozpakowane.
		\item \textbf{Przeglądanie plików}: Lista Twoich plików jest widoczna na stronie głównej. Możesz sortować je po dacie lub nazwie.
		\item \textbf{Pobieranie}: Kliknij nazwę pliku, aby pobrać go na dysk.
		\item \textbf{Uzyskanie opisu AI}: Przy plikach tekstowych (.txt) widoczny jest przycisk "AI Opis". Kliknij go, aby wygenerować podsumowanie. Po wygenerowaniu, przycisk zmieni się na "Pokaż Opis".
		\item \textbf{Usuwanie}: Aby usunąć plik, kliknij ikonę kosza. Możesz też zaznaczyć checkboxami wiele plików i użyć przycisku "Usuń zaznaczone" na dole ekranu.
	\end{itemize}
	
	\newpage
	\section{Wnioski}
	Realizacja projektu chmurowego \textit{GutCloud} dostarczyła zespołowi wielu cennych spostrzeżeń. Przejście przez pełny cykl życia oprogramowania – od wstępnej koncepcji, przez dobór technologii, aż po implementację i wdrożenie – pozwoliło na pełne zrozumienie wyzwań, z jakimi mierzą się współcześni inżynierowie oprogramowania.
	
	Pierwszym kluczowym wnioskiem jest potwierdzenie ogromnej elastyczności i potęgi modelu \textbf{Platform as a Service (PaaS)} oraz \textbf{Infrastructure as a Service (IaaS)}. Wykorzystanie usług zarządzanych (Managed Services), takich jak Azure Blob Storage czy Table Storage, pozwoliło nam zaoszczędzić dziesiątki godzin pracy, które musielibyśmy poświęcić na konfigurowanie własnych serwerów plików, baz danych i systemów backupu. Chmura rzeczywiście demokratyzuje dostęp do technologii enterprise – jako studenci mieliśmy dostęp do tej samej infrastruktury, z której korzystają największe korporacje świata.
	
	W kontekście technologicznym, projekt udowodnił, że język \textbf{Python w połączeniu z frameworkiem Flask} jest doskonałym wyborem do szybkiego prototypowania i budowania aplikacji chmurowych (Rapid Application Development). Biblioteki dostarczone przez Microsoft (Azure SDK for Python) są intuicyjne i dobrze udokumentowane, co pozwoliło na bezproblemową integrację z chmurą. Jednakże, nie obyło się bez wyzwań. Zarządzanie operacjami długotrwałymi, takimi jak rozpakowywanie dużych archiwów ZIP, wymagało zastosowania programowania współbieżnego (wątki). Uświadomiło nam to, że w środowisku webowym nie można blokować głównego wątku obsługi żądań i konieczne jest stosowanie patternów takich jak "Background Workers" czy kolejki zadań.
	
	Integracja z modelem językowym (AI) była najbardziej innowacyjnym elementem projektu. Pokazała ona, że współczesne aplikacje chmurowe ewoluują w stronę systemów inteligentnych. Możliwość automatycznego streszczania dokumentów to funkcja, która realnie zwiększa produktywność użytkownika. Wniosek płynący z tego doświadczenia jest taki, że API modeli generatywnych stają się nowym standardem w toolboxie programisty, podobnie jak kiedyś stały się bazy danych SQL.
	
	Podsumowując, projekt zakończył się sukcesem. Wszystkie założone funkcjonalności zostały zaimplementowane i przetestowane. Aplikacja jest stabilna, skalowalna i bezpieczna. W przyszłości system można by rozwinąć o dodatkowe funkcje, takie jak uwierzytelnianie wieloskładnikowe (MFA), pełnotekstowe wyszukiwanie w dokumentach (Azure Cognitive Search) czy konteneryzację aplikacji przy użyciu Dockera i Kubernetes, co jeszcze bardziej zbliżyłoby ją do komercyjnych standardów produkcyjnych. Projekt ten był doskonałą lekcją inżynierii chmurowej i przygotował nas do pracy w nowoczesnym środowisku IT.
	
	\newpage
	\section{Bibliografia}
	\begin{enumerate}
		\item Dokumentacja Microsoft Azure, \url{https://learn.microsoft.com/en-us/azure/} - dostęp: styczeń 2026.
		\item Dokumentacja Frameworka Flask, \url{https://flask.palletsprojects.com/} - dostęp: styczeń 2026.
		\item Dokumentacja Azure SDK for Python (Storage Blob), \url{https://pypi.org/project/azure-storage-blob/} - dostęp: styczeń 2026.
		\item Groq API Documentation, \url{https://console.groq.com/docs/quickstart} - dostęp: styczeń 2026.
		\item "Cloud Computing Concepts, Technology \& Architecture", Thomas Erl, Ricardo Puttini, Zaigham Mahmood.
		\item Materiały dydaktyczne z przedmiotu "Systemy Przetwarzania w Chmurze".
	\end{enumerate}
	
	\section{Członkowie zespołu}
	\begin{table}[h!]
		\centering
		\renewcommand{\arraystretch}{1.5}
		\begin{tabular}{|p{5cm}|p{9cm}|}
			\hline
			\textbf{Imię i Nazwisko} & \textbf{Rola i Obowiązki} \\
			\hline
			& \textbf{Team Leader / Backend Developer} \newline Projektowanie architektury systemu, konfiguracja usług Azure, implementacja logiki biznesowej w Pythonie, koordynacja prac zespołu. \\
			\hline
			& \textbf{Frontend Developer / UI/UX Designer} \newline Tworzenie szablonów HTML/CSS, projektowanie interfejsu użytkownika, dbałość o responsywność i User Experience. \\
			\hline
			& \textbf{Cloud Engineer / Tester} \newline Zarządzanie subskrypcją Azure, testowanie funkcjonalności, przygotowanie dokumentacji testowej, weryfikacja bezpieczeństwa. \\
			\hline
		\end{tabular}
		\caption{Podział ról i obowiązków w zespole projektowym}
	\end{table}
	
\end{document}